\subsection{Picard 群}\label{src:9.4.5-subsection}
我们介绍 Drezet 与 Narasimhan 对模空间 $M_{r,d}^{\mathrm{ss}}$ 和 $M_{r,\xi}^{\mathrm{ss}}$ 的 Picard 群的计算 \cite[定理 B 和 C]{DN89}. 该论证也将给出模叠 $\cBun_{r,d}^{\mathrm{ss}}$ 与 $\cBun_{r,\xi}^{\mathrm{ss}}$ 的 Picard 群.

秩 $r=1$ 时, 我们实际上已经计算了 Picard 群. Picard 叠到 Picard 概形的带 $\mathbb G_m$-带的 gerbe $\underline{\mathcal P\!ic}^d\to\underline{\mathrm{Pic}}^d$ 平凡 (命题~\bookref{9.4.30}), 即 $\underline{\mathcal P\!ic}^d\cong\underline{\mathrm{Pic}}^d\times B\mathbb G_m$. 因而
\[
\mathrm{Pic}(\underline{\mathcal P\!ic}^d)\cong\mathrm{Pic}(\underline{\mathrm{Pic}}^d)\oplus\mathbb Z.
\]
此外, $A:=\underline{\mathrm{Pic}}^d\cong\underline{\mathrm{Pic}}^0$ 是 Abel 簇, 存在非典范分裂的短正合列
\[
0\longrightarrow\mathrm{Pic}^0(A)\longrightarrow\mathrm{Pic}(A)\longrightarrow\mathrm{NS}(A)\longrightarrow0.
\]
这里 $\mathrm{Pic}^0(A)=A^\vee(\Bbbk)$ 是对偶 Abel 簇, N\'eron--Severi 群是有限生成自由 Abel 群 $\mathbb Z^{\oplus\rho}$. 因此 $\mathrm{Pic}(\underline{\mathrm{Pic}}^d)\cong(\underline{\mathrm{Pic}}^0)^\vee(\Bbbk)\oplus\mathbb Z^{\oplus\rho}$.
\begin{theorem}[\texorpdfstring{$\cBun_{r,d}^{\mathrm{ss}}$ 与 $\cBun_{r,\xi}^{\mathrm{ss}}$ 的 Picard 群}{半稳定向量丛叠的 Picard 群}]\label{src:9.4.36}
设 $C$ 是代数闭域 $\Bbbk$ 上亏格 $g\geq2$ 的光滑连通射影曲线. 设 $r>0$ 和 $d$ 为整数, $\xi\in\mathrm{Pic}^d(C)$, 并置 $r'=r/\gcd(r,d)$, $d'=d/\gcd(r,d)$. 有同构
\[
\begin{aligned}
\mathrm{Pic}(\cBun_{r,\xi}^{\mathrm{ss}})&\cong\mathrm{Pic}(\cBun_{r,\xi})\cong\mathbb Z\langle\mathcal L_V^{1/r'}\rangle\cong\mathbb Z\langle\mathcal L_W\rangle,\\
\mathrm{Pic}(M_{r,\xi}^{\mathrm{ps}})&\cong\mathbb Z\langle L_V\rangle,\\
\mathrm{Pic}(\cBun_{r,d}^{\mathrm{ss}})&\cong\mathrm{Pic}(\cBun_{r,d})\cong\mathrm{Pic}(\underline{\mathcal P\!ic}^d)\oplus\mathbb Z\langle\mathcal L_W\rangle,\\
\mathrm{Pic}(M_{r,d}^{\mathrm{ps}})&\cong\mathrm{Pic}(\underline{\mathrm{Pic}}^d)\oplus\mathbb Z\langle L_V\rangle,
\end{aligned}
\]
其中 $V$ 为 $C$ 上秩 $r'$, 次数 $r'(g-1)-d'$ 的向量丛, $W$ 为 $C$ 上任意次数的线丛, $\mathcal L_V$ 是行列式线丛, $L_V$ 是 $M_{r,d}^{\mathrm{ps}}$ 上拉回为 $\mathcal L_V$ 的线丛. 若 $g\neq2$ 或 $r\neq2$, 上述群还分别与 $\cBun_{r,d}^{\mathrm s}$, $M_{r,d}^{\mathrm s}$, $\cBun_{r,\xi}^{\mathrm s}$ 和 $M_{r,\xi}^{\mathrm s}$ 的 Picard 群等同.
\end{theorem}
\begin{remark}\label{src:9.4.37}
对 $C$ 上向量丛 $Q$, 根据习题~\bookref{9.4.17}, 定义~\bookref{9.4.15} 中 $\cBun_{r,d}^{\mathrm{ss}}$ (分别地, $\cBun_{r,\xi}^{\mathrm{ss}}$) 上的行列式线丛 $\mathcal L_Q$ 仅依赖于 $\mathrm{rk}\,Q$ 与 $\det Q$ (分别地, $\mathrm{rk}\,Q$ 与 $\deg Q$). 但由本定理的证明可知, $\cBun_{r,\xi}^{\mathrm{ss}}$ 上的 $\mathcal L_Q$ 仅依赖于 $\mathrm{rk}\,Q$ (习题~\bookref{9.4.42}). 对线丛 $W$, 在 $\cBun_{r,\xi}^{\mathrm{ss}}$ 上有 $\mathcal L_W^{\otimes r'}\cong\mathcal L_V$, 而 $\mathcal L_W^{\otimes r'}\in\mathrm{Pic}(\underline{\mathcal P\!ic}^d)\oplus\mathbb Z\langle\mathcal L_V\rangle$.
\end{remark}
还应注意, 根据命题~\bookref{9.4.29}, $\mathcal L_V$ 下降为 $M_{r,d}^{\mathrm{ps}}$ 上唯一的线丛 $L_V$.

证明将采用与 $\cBun_{r,\xi}^{\mathrm{ss}}$ 的单有理性 (\bookref{9.4.2}) 的证明相似的策略: 当 $d\gg0$ 时, 每个秩 $r$, 行列式 $\xi\in\mathrm{Pic}^d(C)$ 的半稳定向量丛 $F$ 都可表示为 $\xi$ 被 $\mathcal O_C^{\oplus r-1}$ 的扩张 (引理~\bookref{9.4.3}). 置 $E:=\mathrm{Ext}^1_{\mathcal O_C}(\xi,\mathcal O_C^{\oplus r-1})^\vee$, $C\times\mathbb AE$ 上的万有扩张定义映射 $\mathbb AE\to\cBun_{r,\xi}$, 其像包含半稳定轨迹. 此态射不能下降为映射 $\mathbb PE\to\cBun_{r,\xi}$, 但取 $\mathcal O(r-1)$ 的 $r$ 次根 gerbe $\mathbb PE(\sqrt[r]{\mathcal O(r-1)})$ 后可以下降. 为将 $\mathrm{Pic}(\cBun_{r,\xi})$ 与 $\mathrm{Pic}(\mathbb PE(\sqrt[r]{\mathcal O(r-1)}))\cong\mathbb Z$ 联系起来, 我们将如下分解该映射: $\xi$ 被 $\mathcal O_C^{\oplus r-1}$ 的扩张由一个 $(r-1)$ 维子空间 $W\subseteq\mathrm H^0(C,F)$ 及一个精确到数乘的同构 $\Bbbk^{\oplus r-1}\xrightarrow{\sim}W$ 确定. 换言之, 该根 gerbe 可与 $\cBun_{r,\xi}^{\mathrm{ss}}$ 上相对 Grassmann 簇之上的射影标架丛等同.
\begin{proof}
我们处理固定行列式的情形, 将 $\cBun_{r,d}^{\mathrm{ss}}$ 的情形留给读者 (习题~\bookref{9.4.41}). 对 $p\in C(\Bbbk)$, $n\in\mathbb Z$, 张量 $\mathcal O_C(np)$ 给出同构 $\cBun_{r,\xi}^{\mathrm{ss}}\to\cBun_{r,\xi(rnp)}^{\mathrm{ss}}$, 故可假设 $d=\deg\xi>r(2g-1)$. 置 $E:=\mathrm{Ext}^1_{\mathcal O_C}(\xi,\mathcal O_C^{\oplus r-1})^\vee$. 习题~\bookref{9.2.12}(c) 给出 $C\times\mathbb PE$ 上的万有扩张:
\begin{equation}\label{src:9.4.38}
0\longrightarrow\mathcal O_{C\times\mathbb PE}^{\oplus r-1}\otimes p_2^*\mathcal O_{\mathbb PE}(1)\longrightarrow\mathcal E_{\mathrm{univ}}\longrightarrow p_1^*\xi\longrightarrow0.
\end{equation}
因此 $\det\mathcal E_{\mathrm{univ}}\cong p_1^*\xi\otimes p_2^*\mathcal O_{\mathbb PE}(r-1)$. 记住, $\cBun_{r,\xi}$ 参数化向量丛 $F$ 与同构 $\alpha:\det F\xrightarrow{\sim}\xi$ 所成的对 $(F,\alpha)$. 设 $q:\mathbb PE(\sqrt[r]{\mathcal O(r-1)})\to\mathbb PE$ 是 $\mathcal O(r-1)$ 的 $r$ 次根 gerbe (见例~\bookref{4.9.21} 和习题~\bookref{7.4.26}), $\gamma$ 是该 gerbe 上满足 $\gamma^{\otimes r}\cong(\mathrm{id}\times q)^*\mathcal O_{\mathbb PE}(r-1)$ 的线丛. 则 $C\times\mathbb PE(\sqrt[r]{\mathcal O(r-1)})$ 上的向量丛 $(\mathrm{id}\times q)^*\mathcal E_{\mathrm{univ}}\otimes p_2^*\gamma^\vee$ 及同构
\[
\det\big((\mathrm{id}\times q)^*\mathcal E_{\mathrm{univ}}\otimes p_2^*\gamma^\vee\big)\cong p_1^*\xi
\]
定义映射 $\mathbb PE(\sqrt[r]{\mathcal O(r-1)})\to\cBun_{r,\xi}$. 根据习题~\bookref{7.1.22},
\[
\mathrm{Pic}\big(\mathbb PE(\sqrt[r]{\mathcal O(r-1)})\big)\cong\mathrm{Pic}(\mathbb PE)[\sqrt[r]{\mathcal O(r-1)}]\cong\mathbb Z.
\]
现在构造交换图
\[
\begin{tikzcd}[ampersand replacement=\&,column sep=small]
\mathbb P\mathrm{Fr}(W_{\mathrm{univ}})\arrow[r,"\sim"]\arrow[d,"h"']\&\mathbb PE(\sqrt[r]{\mathcal O(r-1)})^{\mathrm{ss}}\arrow[r,hook]\arrow[d]\&\mathbb PE(\sqrt[r]{\mathcal O(r-1)})\arrow[d]\\
\Gr^{\mathrm{inj}}\arrow[r,"g"']\&\cBun_{r,\xi}^{\mathrm{ss}}\arrow[r,hook]\arrow[ur,phantom,"\square"]\&\cBun_{r,\xi},
\end{tikzcd}
\]
其中左侧的空间将在以下两段中定义.

由于 $d>r(2g-1)$, 每个秩 $r$, 次数 $d$ 的半稳定向量丛均由整体截面生成且无高次上同调 (引理~\bookref{9.2.7}). 设 $\mathcal F_{\mathrm{univ}}$ 为 $C\times\cBun_{r,\xi}^{\mathrm{ss}}$ 上万有向量丛. 上同调与基变换 (\bookref{A.6.8}) 表明 $p_{2,*}\mathcal F_{\mathrm{univ}}$ 是秩 $R:=d+r(1-g)$ 的向量丛, 且 $p_2^*p_{2,*}\mathcal F_{\mathrm{univ}}\to\mathcal F_{\mathrm{univ}}$ 满射. 考虑相对 Grassmann 簇 $g:\Gr(K,p_{2,*}\mathcal F_{\mathrm{univ}})\to\cBun_{r,\xi}^{\mathrm{ss}}$, 参数化维数 $K=R-(r-1)$ 的商 (或维数 $r-1$ 的子空间). 记 $W_{\mathrm{univ}}\subseteq g^*p_{2,*}\mathcal F_{\mathrm{univ}}$ 为万有子丛. 考虑 Cartesian 图
\[
\begin{tikzcd}[ampersand replacement=\&]
C\times\Gr(K,p_{2,*}\mathcal F_{\mathrm{univ}})\arrow[r,"\mathrm{id}\times g"]\arrow[d,"p_2"']\& C\times\cBun_{r,\xi}^{\mathrm{ss}}\arrow[d,"p_2"]\\
\Gr(K,p_{2,*}\mathcal F_{\mathrm{univ}})\arrow[r,"g"']\&\cBun_{r,\xi}^{\mathrm{ss}}\arrow[ul,phantom,"\square"],
\end{tikzcd}
\]
定义 $\Gr^{\mathrm{inj}}\subseteq\Gr(K,p_{2,*}\mathcal F_{\mathrm{univ}})$ 为下列映射的余核是线丛的开轨迹:
\[
p_2^*W_{\mathrm{univ}}\hookrightarrow p_2^*g^*p_{2,*}\mathcal F_{\mathrm{univ}}\cong p_2^*p_{2,*}(\mathrm{id}\times g)^*\mathcal F_{\mathrm{univ}}\twoheadrightarrow(\mathrm{id}\times g)^*\mathcal F_{\mathrm{univ}}.
\]
此线丛必同构于 $\xi$ 的拉回. 根据习题~\bookref{9.4.7}, 补集 $\Gr(K,p_{2,*}\mathcal F_{\mathrm{univ}})\setminus\Gr^{\mathrm{inj}}$ 是次数 $d$ 的超曲面. 滥用记号, 仍以 $W_{\mathrm{univ}}$ 表示万有子丛在 $\Gr^{\mathrm{inj}}$ 上的限制. 令 $h:\mathbb P\mathrm{Fr}(W_{\mathrm{univ}})\to\Gr^{\mathrm{inj}}$ 为射影标架丛 (习题~\bookref{B.1.57}). 留给读者验证 $\mathbb P\mathrm{Fr}(W_{\mathrm{univ}})\cong\mathbb PE(\sqrt[r]{\mathcal O(r-1)})^{\mathrm{ss}}$ (习题~\bookref{9.4.40}). 注意 $\mathbb P\mathrm{Fr}(W_{\mathrm{univ}})$ 是
\[
\widetilde{\mathbb P}:=\mathbb P\big(\mathcal H\!om_{\mathcal O_{\Gr^{\mathrm{inj}}}}(\mathcal O_{\Gr^{\mathrm{inj}}}^{\oplus r-1},W_{\mathrm{univ}})^\vee\big)
\]
中诱导映射 $u:\mathcal O_{\widetilde{\mathbb P}}^{\oplus r-1}\to(W_{\mathrm{univ}})_{\widetilde{\mathbb P}}\otimes\mathcal O_{\widetilde{\mathbb P}}(1)$ 为同构的开轨迹, 即 $\det(u)\neq0$ 的轨迹. 这说明补集 $\widetilde{\mathbb P}\setminus\mathbb P\mathrm{Fr}$ 是次数 $r-1$ 的超曲面.

现在可以计算 Picard 群. 因为 $\cBun_{r,\xi}$ 光滑且 $\operatorname{codim}(\cBun_{r,\xi}\setminus\cBun_{r,\xi}^{\mathrm{ss}})\geq2$ (习题~\bookref{9.2.18}), 有 $\mathrm{Pic}(\cBun_{r,\xi}^{\mathrm{ss}})\cong\mathrm{Pic}(\cBun_{r,\xi})$. 同样,
\[
\mathrm{Pic}\big(\mathbb PE(\sqrt[r]{\mathcal O(r-1)})^{\mathrm{ss}}\big)\cong\mathbb PE(\sqrt[r]{\mathcal O(r-1)}).
\]
由于 $\Gr(K,p_{2,*}\mathcal F_{\mathrm{univ}})\to\cBun_{r,\xi}^{\mathrm{ss}}$ 是相对 Grassmann 簇, 有 $\mathrm{Pic}(\Gr)\cong\mathrm{Pic}(\cBun_{r,\xi}^{\mathrm{ss}})\times\mathbb Z$. 切除序列 $\mathbb Z\to\mathrm{Pic}(\Gr)\to\mathrm{Pic}(\Gr^{\mathrm{inj}})\to0$ 右正合. 补集 $\Gr(K,p_{2,*}\mathcal F_{\mathrm{univ}})\setminus\Gr^{\mathrm{inj}}$ 是次数 $d$ 的超曲面 (习题~\bookref{9.4.7}). 其类可写为 $(N,d)\in\mathrm{Pic}(\cBun_{r,\xi}^{\mathrm{ss}})\times\mathbb Z$, 其中 $N$ 为线丛, 故将 $1$ 映到该补集之类的映射 $\mathbb Z\to\mathrm{Pic}(\Gr(K,p_{2,*}\mathcal F_{\mathrm{univ}}))$ 单射. 得到下列实线箭头组成的交换图, 其中行与列正合:
\[
\begin{tikzcd}[ampersand replacement=\&,column sep=small,row sep=small]
\&\&0\arrow[d]\&0\arrow[d,dashed]\&\\
\&\&\mathrm{Pic}(\cBun_{r,\xi}^{\mathrm{ss}})\arrow[d]\arrow[r,equal]\&\mathrm{Pic}(\cBun_{r,\xi}^{\mathrm{ss}})\arrow[d,"p^*"]\&\\
0\arrow[r]\&\mathbb Z\arrow[r]\arrow[d,equal]\&\mathrm{Pic}(\Gr(K,p_{2,*}\mathcal F_{\mathrm{univ}}))\arrow[r]\arrow[d]\&\mathrm{Pic}(\Gr^{\mathrm{inj}})\arrow[r]\arrow[d,dashed]\&0\\
0\arrow[r]\&\mathbb Z\arrow[r,"d"]\&\mathbb Z\arrow[d]\arrow[r,dashed]\&\mathbb Z/d\arrow[d,dashed]\arrow[r,dashed]\&0\\
\&\&0\&0\&
\end{tikzcd}
\]
我们断言 $p^*:\mathrm{Pic}(\cBun_{r,\xi}^{\mathrm{ss}})\to\mathrm{Pic}(\Gr^{\mathrm{inj}})$ 单射: 若 $p^*(L)=0$, 则 $L$ 在 $\Gr(K,p_{2,*}\mathcal F_{\mathrm{univ}})$ 上的拉回属于到 $\mathrm{Pic}(\Gr^{\mathrm{inj}})$ 的映射的核, 所以存在整数 $\nu$, 使得 $(L,0)=(N^{\otimes\nu},\nu d)\in\mathrm{Pic}(\cBun_{r,\xi}^{\mathrm{ss}})\times\mathbb Z$, 从而 $\nu=0$, $L=0$. 蛇引理表明 $p^*$ 的余核为 $\mathbb Z/d\mathbb Z$, 给出短正合列
\begin{equation*}\tag{$\star$}\label{src:9.4-picard-star}
0\longrightarrow\mathrm{Pic}(\cBun_{r,\xi}^{\mathrm{ss}})\longrightarrow\mathrm{Pic}(\Gr^{\mathrm{inj}})\longrightarrow\mathbb Z/d\longrightarrow0.
\end{equation*}
同样, 因为 $\widetilde{\mathbb P}\to\Gr^{\mathrm{inj}}$ 是相对射影空间, 且 $\mathbb P\mathrm{Fr}(W_{\mathrm{univ}})\subseteq\widetilde{\mathbb P}$ 的补集为次数 $r-1$ 的超曲面, 相同论证给出短正合列
\begin{equation*}\tag{$\star\star$}\label{src:9.4-picard-double-star}
0\longrightarrow\mathrm{Pic}(\Gr^{\mathrm{inj}})\longrightarrow\mathrm{Pic}(\mathbb P\mathrm{Fr}(W_{\mathrm{univ}}))\longrightarrow\mathbb Z/(r-1)\longrightarrow0.
\end{equation*}
序列~\eqref{src:9.4-picard-star} 和~\eqref{src:9.4-picard-double-star} 合起来表明, $\mathrm{Pic}(\cBun_{r,\xi}^{\mathrm{ss}})\subseteq\mathrm{Pic}(\mathbb P\mathrm{Fr}(W_{\mathrm{univ}}))\cong\mathbb Z$ 的余核为 $\mathbb Z/d(r-1)$. 这已经说明 $\mathrm{Pic}(\cBun_{r,\xi}^{\mathrm{ss}})\cong\mathbb Z$.

为确定一个生成元, 对 $C$ 上任意向量丛 $Q$, 计算行列式线丛 $\mathcal L_Q$ 到 $\mathbb PE(\sqrt[r]{\mathcal O(r-1)})$ 的拉回. 将定义万有扩张的序列~\bookeqref{9.4.38} 沿 $\mathrm{id}\times q:C\times\mathbb PE\to C\times\mathbb PE(\sqrt[r]{\mathcal O(r-1)})$ 拉回, 并张量 $p_2^*\gamma^\vee$ (回忆 $\gamma^{\otimes r}\cong(\mathrm{id}\times q)^*\mathcal O_{\mathbb PE}(r-1)$), 得到正合列
\[
\begin{aligned}
0&\longrightarrow\mathcal O_{C\times\mathbb PE(\sqrt[r]{\mathcal O(r-1)})}^{\oplus r-1}\otimes p_2^*(q^*\mathcal O_{\mathbb PE}(1)\otimes\gamma^\vee)\\
&\longrightarrow\underbrace{(\mathrm{id}\times q)^*\mathcal E_{\mathrm{univ}}\otimes p_2^*\gamma^\vee}_{\mathcal F_{\mathrm{univ}}|_{C\times\mathbb PE(\sqrt[r]{\mathcal O(r-1)})}}\longrightarrow p_1^*\xi\otimes p_2^*\gamma^\vee\longrightarrow0.
\end{aligned}
\]
再张量 $p_1^*Q$ 并应用 $\mathrm Rp_{2,*}$, 得到根 gerbe 上的正合三角
\[
\begin{aligned}
\mathrm R\Gamma(C,Q)^{\oplus r-1}\otimes(q^*\mathcal O_{\mathbb PE}(1)\otimes\gamma^\vee)&\longrightarrow\mathrm Rp_{2,*}(\mathcal F_{\mathrm{univ}}\otimes p_1^*Q)|_{\mathbb PE(\sqrt[r]{\mathcal O(r-1)})}\\
&\longrightarrow\mathrm R\Gamma(C,Q\otimes\xi)\otimes\gamma^\vee.
\end{aligned}
\]
取行列式得到
\[
\begin{aligned}
\mathcal L_Q|_{\mathbb PE(\sqrt[r]{\mathcal O(r-1)})}&=(r-1)\chi(C,Q)(\gamma-q^*\mathcal O_{\mathbb PE}(1))+\chi(C,Q\otimes\xi)\gamma\\
&=d\,\mathrm{rk}(Q)\gamma.
\end{aligned}
\]
$\mathrm{Pic}(\mathbb PE(\sqrt[r]{\mathcal O(r-1)}))$ 的一个生成元可写为 $G=q^*\mathcal O_{\mathbb PE}(1)-\gamma$, 从而 $q^*\mathcal O_{\mathbb PE}(1)=rG$, $\gamma=(r-1)G$. 取 $W$ 为秩 $1$ 的向量丛, 得 $\mathcal L_W|_{\mathbb PE(\sqrt[r]{\mathcal O(r-1)})}=d\gamma=d(r-1)G$, 因此 $\mathcal L_W$ 自由生成 $\mathrm{Pic}(\cBun_{r,\xi}^{\mathrm{ss}})$. 若 $V$ 是秩 $r'$, 次数 $r'(g-1)-d'$ 的向量丛, 则 $\mathcal L_V|_{\mathbb PE(\sqrt[r]{\mathcal O(r-1)})}=r'd(r-1)G$, 且 $\mathcal L_V=\mathcal L_W^{\otimes r'}$.

最后一个断言成立, 因为当 $g\neq2$ 或 $r\neq2$ 时, 严格半稳定轨迹的余维至少为二 (习题~\bookref{9.2.18}).
\end{proof}
\begin{remark}\label{src:9.4.39}
互素情形, 即 $\gcd(r,d)=1$, 由 Seshadri 给出; 见 \cite[命题 3.4]{NR69}. $r=2$ 的情形由 Beauville 给出 \cite[定理 1]{Bea88}. 一般情形见 \cite[定理 B 和 C]{DN89}.
\end{remark}
\begin{exercise}\label{src:9.4.40}
利用在 $\cBun_{r,\xi}^{\mathrm{ss}}$ 上的函子描述, 将 $\mathbb P\mathrm{Fr}(W_{\mathrm{univ}})$ 与 $\mathbb PE(\sqrt[r]{\mathcal O(r-1)})^{\mathrm{ss}}$ 等同.
\end{exercise}
\begin{exercise}\label{src:9.4.41}
沿用上述策略, 但在 $\underline{\mathcal P\!ic}^d$ 上相对地进行, 计算 $\cBun_{r,d}^{\mathrm{ss}}$ 的 Picard 群.

\textit{提示:} 设 $\xi_{\mathrm{univ}}$ 是 $C\times\underline{\mathcal P\!ic}^d$ 上的万有线丛, 置
\[
\mathbb E:=\big(p_{2,*}\mathcal E\!xt^1_{\mathcal O_{C\times\underline{\mathcal P\!ic}^d}}(\xi_{\mathrm{univ}},\mathcal O_{C\times\underline{\mathcal P\!ic}^d})\big)^\vee,
\]
并考虑 $\underline{\mathcal P\!ic}^d$ 上的态射 $\mathbb P\mathbb E(\sqrt[r]{\mathcal O(r-1)})\to\cBun_{r,d}^{\mathrm{ss}}$.
\end{exercise}
\begin{exercise}\label{src:9.4.42}
证明 $\cBun_{r,d}^{\mathrm{ss}}$ 上的 $\mathcal L_Q\otimes\det^*\mathcal P_Q^\vee$ (分别地, $\cBun_{r,\xi}^{\mathrm{ss}}$ 上的 $\mathcal L_Q$) 仅依赖于 $\mathrm{rk}\,Q$, 其中 $\mathcal P_Q$ 为 $\underline{\mathcal P\!ic}^d$ 上的行列式线丛, $\det:\cBun_{r,d}^{\mathrm{ss}}\to\underline{\mathcal P\!ic}^d$ 是行列式态射.

\textit{提示:} 使用 $\mathrm{Pic}(\cBun_{r,d}^{\mathrm{ss}})\to\mathrm{Pic}(\mathbb P\mathbb E)$ 的单射性及 $\mathcal L_Q|_{\mathbb P\mathbb E}$ 的公式.
\end{exercise}
\begin{corollary}\label{src:9.4.43}
模空间 $M_{r,\xi}^{\mathrm{ps}}$ 是 Fano 簇.
\end{corollary}
\begin{proof}
证明中假设 $M_{r,\xi}^{\mathrm{ps}}$ 射影 (稍后将在 \S\bookref{9.5}--\bookref{9.6} 中以两种方式证明). 因为 $M_{r,\xi}^{\mathrm{ps}}$ 单有理 (\bookref{9.4.2}), 典范除子 $K_{M_{r,\xi}^{\mathrm{ps}}}$ 既非丰沛也非平凡. 由于 $\mathrm{Pic}(M_{r,\xi}^{\mathrm{ps}})\cong\mathbb Z$, $-K_{M_{r,\xi}^{\mathrm{ps}}}$ 必为丰沛的.
\end{proof}
\begin{exercise}\label{src:9.4.44}
将 $M_{r,\xi}^{\mathrm{ss}}$ 的对偶化层与 $L_V^{\otimes-2k}$ 等同, 其中 $k=\gcd(r,\deg\xi)$. 另见 \cite[定理 F]{DN89}.
\end{exercise}
\begin{exercise}[难]\label{src:9.4.45}
证明 $M_{r,\xi}^{\mathrm{ss}}$ 的 Brauer 群为 $\mathbb Z/\gcd(r,\deg\xi)$.

\textit{提示:} 沿用定理~\bookref{9.4.36} 的策略. 另见 \cite[定理 1.8]{BBGN07}.
\end{exercise}
